% Compile from the Time series/ folder:
%   pdflatex -shell-escape report.tex
%   bibtex report
%   pdflatex -shell-escape report.tex
%   pdflatex -shell-escape report.tex

\documentclass{article}
\usepackage{microtype}
\usepackage{graphicx}
\usepackage{booktabs}
\usepackage{hyperref}
\usepackage{amsmath,amssymb}
\usepackage{xcolor}
\usepackage[accepted]{ICML26/icml2026}

\icmltitlerunning{Foundation Model Adaptation for LSST Time Series Classification}

\begin{document}

\twocolumn[
  \icmltitle{Adapting Foundation Models for Astronomical Time Series Classification:\\
             Setting 1 on the LSST/PLAsTiCC Benchmark}

  \begin{icmlauthorlist}
    \icmlauthor{Zi\'e Coulibaly}{u1}
    \icmlauthor{Clijo Jose}{u1}
    \icmlauthor{Oumouhani El~Vilaly}{u1}
  \end{icmlauthorlist}

  \icmlaffiliation{u1}{Deep Learning for Time Series, 2026}

  \vskip 0.3in
]

\printAffiliationsAndNotice{}

\begin{abstract}
We address 14-class astronomical object classification on the LSST/PLAsTiCC
benchmark~\citep{plasticc} — a multivariate time series dataset of shape
$(N,\,T{=}36,\,C{=}6)$ with severe class imbalance (155:1 ratio) and sparse
observations. Following Setting~1, we adapt two pre-trained foundation models:
MOMENT-1-large \citep{moment} and Chronos-T5-Small \citep{chronos}. We further
compare with UniTS \citep{units}, a non-pretrained dual-attention architecture
trained with data augmentation. Our best ensemble — combining Chronos,
InceptionTime$\times$5, PatchTST, and MultiROCKET via F1-weighted soft voting —
achieves \textbf{64.48\% accuracy} and \textbf{Weighted F1 = 0.61}, comparing
favourably to the best published result on this benchmark (MUSE, 63.6\%, \citealt{bakeoff}).
Code: \url{https://github.com/coulzie031/DeepLearningForTimeSeries}
\end{abstract}

% ─────────────────────────────────────────────────
\section{Introduction}
% ─────────────────────────────────────────────────

The LSST (Large Synoptic Survey Telescope) dataset, derived from the PLAsTiCC
challenge, is one of the most challenging multivariate time series classification
(TSC) benchmarks. Each light curve records the brightness of an astronomical
object across 36 epochs and 6 photometric filters (passbands $u,g,r,i,z,y$),
yielding a tensor of shape $(N,36,6)$. Two factors make this dataset
particularly hard: (i) \emph{sparse observations} leading to NaN values, and
(ii) \emph{extreme class imbalance} — class 53 (kilonovae) has only 7 training
samples out of 2,459 total.

Foundation models pre-trained on large time-series corpora offer a natural
approach to scarce-data regimes. We explore two such models under
\emph{Setting~1}: adapt a pre-trained foundation model to the classification
task. We further investigate UniTS \citep{units}, a from-scratch dual-attention
architecture combined with targeted data augmentation, as a strong competitor.
Finally, we combine complementary models through an ensemble to push accuracy
beyond the published SOTA.

% ─────────────────────────────────────────────────
\section{Dataset and Preprocessing}
% ─────────────────────────────────────────────────

\paragraph{LSST/PLAsTiCC.}
The dataset contains 2,459 training samples, 2,466 test samples, and 14 classes
(supernovae subtypes, variable stars, AGN, kilonovae, microlensing events).
Labels are heavily imbalanced: the dominant class (SN~Ia, class~90) represents
44\% of training data while class~53 has 7 samples.

\paragraph{Preprocessing.}
Missing values are filled via forward-fill, backward-fill, then zero-fill. A
binary mask $(B,T,C)$ encodes observation positions. All channels are
z-normalized using training statistics. To address imbalance, we use a
\texttt{WeightedRandomSampler} with inverse-frequency weights and
Focal Loss ($\gamma=2$) \citep{focal}, which downweights easy majority examples.

\paragraph{Data augmentation} (applied only during training):
\emph{Jitter} adds Gaussian noise ($\sigma{=}0.05$) with 50\% probability;
\emph{scaling} multiplies each channel by a random factor in $[0.9, 1.1]$;
\emph{time shift} applies a circular roll of $\pm 3$ steps with 30\%
probability; \emph{MixUp} ($\alpha{=}0.3$) interpolates pairs of samples and
their labels.

% ─────────────────────────────────────────────────
\section{Methods}
% ─────────────────────────────────────────────────

\subsection{Baseline: InceptionTime}

InceptionTime \citep{inception} applies multi-scale 1D convolutions to each
channel independently. We use kernel sizes $(9,19,39)$, 2 residual Inception
blocks, $F{=}32$ filters, and Global Average Pooling — totalling 474k
parameters. No pre-training is used. Trained with Focal Loss, AdamW, and
cosine annealing, early-stopped on validation macro F1.

\subsection{Setting 1 — Foundation Model Adaptation}

\paragraph{MOMENT-1-large \citep{moment}.}
A 341M-parameter T5-based encoder pre-trained on over 1B timesteps from
diverse domains. We pad sequences to $T{=}40$ (multiple of patch size~8),
extract per-channel embeddings of dimension 1024, and flatten to a 6144-dim
vector fed to an MLP head. Training proceeds in two phases: (1) linear
probing — backbone frozen, only the head trains (200 epochs, early stop
ep62, val F1 = 0.303); (2) partial unfreeze — last 4 of 24 T5 encoder blocks
unfrozen with $\text{LR}{=}5{\times}10^{-6}$ (100 epochs, early stop ep26,
val F1 = 0.303). Despite native support for classification, MOMENT fails to
improve: 341M parameters is an ill-matched scale for 1,967 training samples.

\paragraph{Chronos-T5-Small \citep{chronos}.}
A 46M-parameter T5-based model originally pre-trained for univariate
probabilistic \emph{forecasting} on real and synthetic time series. We discard
the decoder and repurpose the encoder for classification — analogous to using
BERT for sentence classification. Each of the 6 passbands is tokenized
independently; masked mean-pooling over hidden states yields a 512-dim
per-channel embedding; the 6 embeddings are concatenated into a 3072-dim
representation, which is fed to a 3-layer MLP head.
Training: phase~1 linear probing (200 ep, stop ep86, val F1 = 0.433);
phase~2 partial unfreeze of last 4 encoder blocks,
$\text{LR}_\text{head}{=}5{\times}10^{-5}$,
$\text{LR}_\text{enc}{=}5{\times}10^{-6}$ (100 ep, stop ep32, val F1 = 0.433).
Chronos outperforms MOMENT decisively despite being designed for forecasting,
confirming that model scale matters more than task alignment when data is scarce.

\subsection{Competitor: UniTS with Data Augmentation}

We also evaluate a complementary augmentation-based approach (external
repo$^\dagger$) applying jitter ($\sigma{=}0.05$), channel scaling (${\times}[0.9,1.1]$),
and time shift (${\pm}3$ steps) during training.
This augmentation improves InceptionTime from acc=0.5483/F1=0.3753 to 0.5783/0.4324,
and is used to train PatchTST (0.5985/0.4599) and UniTS (0.6606/0.5314).

UniTS \citep{units} (NeurIPS 2024) interleaves per block: sequence attention
(over time), variable attention (across channels), and a Dynamic Linear
Operator (DLO) for low-rank temporal mixing $W{=}UV^\top$.
Config: $d{=}128$, 8 heads, 6 blocks, rank $r{=}16$ (1.6M params).
Achieves acc = 0.6606, macro F1 = 0.5314 — the highest single-model scores.

\subsection{Additional Models}

\textbf{InceptionTime-Large $\times$5:} 5 seeds, $F{=}64$, 3 blocks + TTA ($N{=}5$). Acc = 0.6427.

\textbf{PatchTST \citep{patchtst}:} channel-independent patches (len=4), $d{=}64$, 3 layers + TTA. Acc = 0.5985, F1 = 0.4599.

\textbf{MultiROCKET \citep{rocket}:} 79,968 random kernel features + RidgeCV. Acc = 0.6079.

\subsection{Ensemble: F1-Weighted Soft Voting}

Each neural model outputs a probability vector $\mathbf{p}_i \in \mathbb{R}^{14}$
over classes. We combine models via weighted soft voting:
\begin{equation}
  \hat{y} = \arg\max_c \sum_i w_i\, p_{i,c},
  \quad
  w_i = \frac{\exp(10 \cdot F1_{\text{val},i})}{\sum_j \exp(10 \cdot F1_{\text{val},j})}
  \label{eq:ensemble}
\end{equation}
Models with validation macro F1 $<$ 0.33 are excluded.
The final ensemble includes InceptionTime-Large$\times$5 ($w{=}0.852$),
PatchTST ($w{=}0.068$), Chronos ($w{=}0.051$), and MultiROCKET ($w{=}0.029$).
MOMENT is excluded (val F1 = 0.303).

% ─────────────────────────────────────────────────
\section{Results}
% ─────────────────────────────────────────────────

Table~\ref{tab:results} summarises all results on the LSST test set (2,466
samples). Macro F1 is low due to two structurally unlearnable classes:
class~53 (7 train samples, F1 = 0) and class~64 (24 samples, F1 = 0.08).
Excluding these two classes, macro F1 rises to $\approx$0.49.
Weighted F1 = 0.61 (weighted by class support) better reflects real-world
performance on the learnable 90\% of test samples.

\begin{table}[t]
\caption{Test set results on LSST (14-class). Acc = accuracy, Mac.F1 = macro
  F1, W.F1 = weighted F1. $^\dagger$External implementation
  (\url{github.com/mhani6/deep\_learning\_time\_series}).
  SOTA from \citet{bakeoff} reports accuracy only.}
\label{tab:results}
\vskip 0.1in
\begin{center}
\begin{small}
\begin{tabular}{lccc}
\toprule
\textbf{Method} & \textbf{Acc} & \textbf{Mac.F1} & \textbf{W.F1} \\
\midrule
\multicolumn{4}{l}{\textit{Baselines (no pre-training)}} \\
InceptionTime                      & 0.5483 & 0.3753 & — \\
InceptionTime + aug.$^\dagger$     & 0.5783 & 0.4324 & — \\
\midrule
\multicolumn{4}{l}{\textit{Foundation models (Setting 1)}} \\
MOMENT-1-large                     & 0.3078 & 0.2791 & — \\
Chronos-T5-Small                   & 0.5333 & 0.4333 & — \\
\midrule
\multicolumn{4}{l}{\textit{Competitor models}} \\
PatchTST + aug.$^\dagger$ + TTA    & 0.5985 & 0.4599 & — \\
MultiROCKET                        & 0.6079 & 0.3625 & — \\
UniTS + aug.$^\dagger$             & 0.6606 & \textbf{0.5314} & — \\
InceptionTime-L$\times$5 + TTA    & 0.6427 & 0.4150 & — \\
\midrule
\textbf{Ensemble (ours)}           & \textbf{0.6448} & 0.4200 & \textbf{0.61} \\
\midrule
\multicolumn{4}{l}{\textit{Published SOTA \citep{bakeoff}}} \\
MUSE                               & 0.636  & —      & — \\
ROCKET                             & 0.632  & —      & — \\
MrSEQL                             & 0.603  & —      & — \\
\bottomrule
\end{tabular}
\end{small}
\end{center}
\vskip -0.1in
\end{table}

\paragraph{Foundation model comparison.}
Chronos (46M) significantly outperforms MOMENT (341M): +22.5 pts accuracy,
+15.4 pts macro F1. The key factor is scale: at 1,967 training samples, a
341M-parameter model is prone to underfitting even under linear probing, while
Chronos' compact encoder fine-tunes effectively. Chronos also surpasses the
baseline (macro F1 0.433 vs.\ 0.375), validating Setting~1.

\paragraph{Competitor comparison.}
UniTS$^\dagger$ achieves the highest single-model scores among all tested methods
(acc = 0.6606, F1 = 0.5314), driven by its dual sequence--variable attention
and low-rank temporal mixing. Our ensemble reaches acc = 0.6448, above the
best published result on this benchmark (MUSE, 63.6\%, \citealt{bakeoff}),
though a direct controlled comparison is limited since \citet{bakeoff} reports
accuracy only and does not address class imbalance.

\paragraph{Class imbalance.}
Both approaches substantially benefit from imbalance handling. Standard
InceptionTime from the bake-off (no imbalance handling) achieves only 34.0\%
\citep{bakeoff}; our baseline with WeightedSampler and Focal Loss reaches
54.8\% — a +20.8 pt gain from preprocessing alone.

% ─────────────────────────────────────────────────
\section{Conclusion}
% ─────────────────────────────────────────────────

We present a comprehensive study of foundation model adaptation (Setting~1) and
complementary architectures on the LSST benchmark. Our key findings are:
(1)~Chronos-T5-Small, despite being designed for forecasting, transfers
effectively to multivariate classification, outperforming MOMENT by a large
margin — confirming that model scale relative to training set size matters more
than pre-training task alignment;
(2)~UniTS with data augmentation achieves the best single-model macro F1
(0.5314), showing the complementary value of dual-attention architectures;
(3)~our ensemble (64.48\%) compares favourably to the best published result
on this benchmark \citep{bakeoff}, without any additional data;
(4)~imbalance handling (WeightedSampler + Focal Loss) is the single most
impactful preprocessing decision on this dataset.

\section*{Impact Statement}
This paper presents work whose goal is to advance the field of Machine Learning.
There are many potential societal consequences of our work, none which we feel
must be specifically highlighted here.

\bibliography{report}
\bibliographystyle{ICML26/icml2026}

\end{document}
